% Copyright (c) 2026 Zixuan Zhang. Distributed under the MIT License.
\documentclass[11pt]{article}
\usepackage[a4paper,margin=30mm]{geometry}
\usepackage{micfong-sans-math}
\usepackage[hidelinks,bookmarks=false,pdfpagelabels=false]{hyperref}

\setmainfont{lmsans10-regular.otf}[
  BoldFont = lmsans10-bold.otf, ItalicFont = lmsans10-oblique.otf]
\setlength{\parindent}{0pt}
\setlength{\parskip}{5pt}
\pagestyle{empty}

\begin{document}

{\LARGE\bfseries Micfong Sans Math}\par
Version 0.10 \hfill Zixuan Zhang

\section*{Use in a document}

Compile with LuaLaTeX or XeLaTeX. Add the following line to your preamble.

\begin{verbatim}
\usepackage{micfong-sans-math}
\end{verbatim}

The package loads \texttt{unicode-math} and selects Micfong Sans Math.
Choose your text font separately with \verb|\setmainfont|. The math font was
designed to pair with IBM Plex Sans and can be used with any text font.

\[
  f(x) = \sum_{n=1}^{\infty} \frac{\sin(nx)}{n^2},
  \qquad \int_0^{2\pi} f(x)^2\,\mathrm{d}x
    = \pi \sum_{n=1}^{\infty} \frac{1}{n^4}.
\]

Use \verb|\symbf| to bold individual symbols. To switch an entire formula to
the bold font, use \verb|\boldmath| outside math mode, as in
\verb|{\boldmath $\alpha + x^2 = 1$}|. The bold font has a smaller character
set than the regular font, including limited coverage of specialised symbols.

\[
  \symbf{A} \symbf{x} = \lambda \symbf{x}
  \qquad
  \text{Regular } \alpha + x^2 = 1
  \qquad
  \text{Bold } \text{\boldmath$\alpha + x^2 = 1$}
\]

\section*{Installation}

Install \texttt{micfong-sans-math} through your TeX distribution's package
manager when available. For a single project, copy the following files from
the archive into the folder that contains your main \texttt{.tex} file.

\begin{verbatim}
micfong-sans-math.sty
MicfongSansMath.otf
MicfongSansMath-Bold.otf
\end{verbatim}

Your TeX distribution must also provide \texttt{unicode-math} and \texttt{iftex}.

\section*{Credits and license}

Micfong Sans Math derives from Lete Sans Math by Chenjing Bu and Daniel Flipo,
with letters and numbers drawn from IBM Plex Sans and IBM Plex Mono.
The fonts are licensed under the SIL Open Font License 1.1. The LaTeX package,
documentation, and build scripts use the MIT License.

\href{https://github.com/micfong-z/micfong-sans-math}{Project repository}
\quad
\href{https://github.com/micfong-z/micfong-sans-math/issues}{Issue tracker}

\end{document}
